% 66e Olympiade Internationale de Mathématiques — Sunshine Coast, Australie
% Énoncés officiels en français (imo-official.org). Chaque problème vaut 7 points.

\begin{center}
  \emph{Premier jour --- mardi 15 juillet 2025 --- Durée : 4 heures 30}
\end{center}

\subsection*{Problème 1}
On dit qu'une droite du plan est \emph{ensoleillée} lorsqu'elle n'est
parallèle ni à l'axe des abscisses, ni à l'axe des ordonnées, ni à la
droite d'équation $x + y = 0$.

Soit $n \geqslant 3$ un entier donné. Trouver tous les entiers
$k \geqslant 0$ pour lesquels il existe $n$ droites du plan, deux à
deux distinctes, telles que
\begin{itemize}
  \item pour tous les entiers $a \geqslant 1$ et $b \geqslant 1$ tels
  que $a + b \leqslant n + 1$, le point de coordonnées $(a, b)$
  appartient à au moins une de ces $n$ droites, et
  \item exactement $k$ de ces $n$ droites sont ensoleillées.
\end{itemize}

\subsection*{Problème 2}
Soit $\Omega$ un cercle de centre $M$ et $\Gamma$ un cercle de centre
$N$ tels que le rayon de $\Omega$ soit strictement plus petit que le
rayon de $\Gamma$. On suppose que les cercles $\Omega$ et $\Gamma$ se
coupent en deux points $A$ et $B$ distincts. La droite $(MN)$ coupe
$\Omega$ en un point $C$ et $\Gamma$ en un point $D$ de sorte que les
points $C$, $M$, $N$ et $D$ soient alignés dans cet ordre. Soit $P$ le
centre du cercle circonscrit au triangle $ACD$. La droite $(AP)$
recoupe $\Omega$ en un point $E$ distinct de $A$ ; elle recoupe
également $\Gamma$ en un point $F$ distinct de $A$. Enfin, soit $H$
l'orthocentre du triangle $PMN$.

Démontrer que la parallèle à $(AP)$ passant par $H$ est tangente au
cercle circonscrit au triangle $BEF$.

(L'\emph{orthocentre} d'un triangle est le point d'intersection de ses
hauteurs.)

\subsection*{Problème 3}
Soit $\N_{\geqslant 1}$ l'ensemble des entiers strictement positifs.
On dit qu'une fonction $f \colon \N_{\geqslant 1} \to \N_{\geqslant 1}$
est \emph{balaise} si
\[
f(a) \ \text{divise} \ b^a - f(b)^{f(a)}
\]
pour tous les entiers $a \geqslant 1$ et $b \geqslant 1$.

Trouver le plus petit nombre réel $c$ tel que $f(n) \leqslant cn$ pour
toute fonction $f$ balaise et tout entier $n \geqslant 1$.

\bigskip
\begin{center}
  \emph{Second jour --- mercredi 16 juillet 2025 --- Durée : 4 heures 30}
\end{center}

\subsection*{Problème 4}
Un \emph{diviseur strict} d'un entier $m \geqslant 1$ est un diviseur
de $m$ strictement positif et différent de $m$.

La suite $a_1, a_2, \ldots$ consiste en une infinité d'entiers
strictement positifs dont chacun possède au moins trois diviseurs
stricts. Pour tout $n \geqslant 1$, l'entier $a_{n+1}$ est la somme
des trois plus grands diviseurs stricts de $a_n$.

Trouver toutes les valeurs possibles de $a_1$.

\subsection*{Problème 5}
Anna et Baptiste jouent au jeu de l'\emph{inékoalaté}, un jeu à deux
joueurs dont les règles dépendent d'un réel $\lambda > 0$ connu des
deux joueurs. Pour tout $n \geqslant 1$, lors du $n^{\text{ème}}$ tour
de jeu,
\begin{itemize}
  \item si $n$ est impair, Anna choisit un réel $x_n \geqslant 0$ tel
  que
  \[
  x_1 + x_2 + \cdots + x_n \leqslant \lambda n ;
  \]
  \item si $n$ est pair, Baptiste choisit un réel $x_n \geqslant 0$
  tel que
  \[
  x_1^2 + x_2^2 + \cdots + x_n^2 \leqslant n.
  \]
\end{itemize}
Si un joueur ne peut pas choisir de réel $x_n$ adéquat, la partie
s'arrête ; il perd et son adversaire gagne. Si la partie ne s'arrête
jamais, aucun joueur ne gagne. À tout moment, chaque joueur connaît la
liste de tous les réels déjà choisis.

Trouver toutes les valeurs de $\lambda$ pour lesquelles Anna dispose
d'une stratégie gagnante, ainsi que toutes celles pour lesquelles
Baptiste dispose d'une stratégie gagnante.

\subsection*{Problème 6}
Clara a dessiné un quadrillage formé de $2025 \times 2025$ carrés
unité. Elle souhaite placer sur ce quadrillage des tuiles
rectangulaires, de tailles possiblement différentes, de sorte que les
côtés de chaque tuile appartiennent à des droites du quadrillage et
que chaque carré unité soit recouvert par au plus une tuile.

Trouver le nombre minimal de tuiles que Clara doit placer pour faire
en sorte que chaque ligne et chaque colonne du quadrillage contiennent
exactement un carré unité qui n'est recouvert par aucune tuile.
